% Owned proposed B.139. All checkpoint28 sources remain frozen.
\subsection*{An actual one-family quotient and its complete fixed-source comparison}

\statement{B.139}{Proved}{The actual one-family quotient has a complete graph presentation and exact discrete and completed fixed-source image; the original comparison reduces to its actual remaining kernel}
\label{stmt:B139}
Use the complete marked group $\Gamma_0$, source operators and selected
involution $\delta=\operatorname{Inn}(\tau)$ of B.123 and B.135.
In the minus local tuple
\[
 E_0=F(a,b,d,e,f)*C_2(C)*C_2(\tau)
\]
fill $b,d,e$ and their selected mates by their full normal closure
in the actual ordinary group. Denote the resulting actual quotient by
$\pi:\Gamma_0\to\Gamma_a$, put $K_{\rm rem}=\ker\pi$, and let
$H_{\rm src}$ be the actual six-source image in $\Gamma_0$.
Every signed source translate and both ordinary node sheets remain
part of the quotient presentation. Write
$J_a=\pi(H_{\rm src})$ and $a^*=\tau a\tau$.
Then
\begin{equation}\label{eq:B139-main}
 \Gamma_a^\delta=J_a=\langle aa^*,m,v\rangle,\qquad
 (\widehat{\Gamma_a})^\delta=\overline{J_a}.
\end{equation}
The second closure is the actual closure in the full completion,
not a presumed abstract completion of $J_a$.

The quotient has an exact graph presentation as follows. Put
\[
 Q_f=C_2(C)*\mathbb Z(f)*C_2(\tau),\qquad S=F(m,v),\qquad Q=Q_f\rtimes S.
\]
The action and its marked conjugator cocycle are specified in the proof.
The stabilizer $H=\langle m,fv\rangle\leq Q$ maps isomorphically onto $S$.
Define a graph $\mathcal G$ with vertices $Q/H$ and edges
\begin{equation}\label{eq:B139-graph}
 \{qH,qCH\},\quad \{qH,q\tau H\}\qquad(q\in Q).
\end{equation}
Repeated edges are identified. If $A(\mathcal G)$ is its graph group,
then the actual orbifold quotient and its ordinary subgroup are
\begin{equation}\label{eq:B139-split}
 E_a=\Gamma_a\rtimes\langle\tau\rangle
       =A(\mathcal G)\rtimes Q,\qquad
 \Gamma_a=A(\mathcal G)\rtimes Q^0 .
\end{equation}
Here the original orientation is odd on both cones and even on $f,S$
and all graph generators; $Q^0$ is its kernel.

For the original group there are the exact, generally nonsplit, reductions
\begin{equation}\label{eq:B139-kernel-reductions}
 \begin{gathered}
 1\longrightarrow K_{\rm rem}^{\delta}
 \longrightarrow\Gamma_0^\delta
 \longrightarrow\Gamma_a^\delta\longrightarrow1,\qquad
 \Gamma_0^\delta=H_{\rm src}K_{\rm rem}^\delta,\\
 (\widehat{\Gamma_0})^\delta
       =\overline{H_{\rm src}}\,(\overline{K_{\rm rem}})^\delta .
 \end{gathered}
\end{equation}
Consequently the original discrete and completed fixed-source equalities
are respectively equivalent to
\begin{equation}\label{eq:B139-remaining-gates}
 \begin{gathered}
 K_{\rm rem}^\delta=H_{\rm src}\cap K_{\rm rem},\\
 (\overline{K_{\rm rem}})^\delta
       =\overline{H_{\rm src}}\cap\overline{K_{\rm rem}}.
 \end{gathered}
\end{equation}
All bars in the last two displays denote actual closures in
$\widehat{\Gamma_0}$. Neither remaining equality is asserted.

\begin{proof}
\emph{B.139.1. The entire actual filling and its surviving relations.}
Use the two complete local bases, actual mirror and both signed
operators printed in B.135. Filling $b,d,e$ is equivalently filling
the neck letters $B=CbC,d,e$. The physical basis of B.135 becomes
\[
 (O_1,O_2,O_3,O_4,O_5)=(1,a,1,1,f).
\]
Its actual mirror formula therefore gives
\[
 (a_+,B_+,d_+,e_+,f_+)=(a^{-1},1,1,1,af^{-1}a^{-1}),
 \qquad C_+=a\tau a^{-1}.
\]
The normal closure of the three deleted neck letters is preserved
by the mirror and its inverse, and by the partial inner operators
in the two local bases. Thus the deleted normal closure is invariant
under both signed source generators. This is a quotient of the
whole actual presentation.

The surviving operators on $\mathbb Z(a)*Q_f$ are
\begin{equation}\label{eq:B139-operators}
 \begin{array}{c|cccc}
       &a&C&f&\tau\\ \hline
 \alpha_m&a&f^{-1}Cf&f^{-1}CfCf&\tau\\
 \alpha_v&f^{-1}af&f^{-1}Cf&\tau f\tau&\tau .
 \end{array}
\end{equation}
The first is partial $\operatorname{Inn}(f^{-1}C)$ on $\langle C,f\rangle$.
Its inverse is partial $\operatorname{Inn}(Cf)$.
For the second, the inverse sends $f$ to $\tau f\tau$, fixes $\tau$,
and conjugates each of $a,C$ by $\tau f\tau$.
These formulas prove both signed actions preserve $Q_f$.

In the ten-word scheme of B.135, only the diagonal-$a$ word
survives in each region. Their complete based images are
\begin{equation}\label{eq:B139-surviving-nodes}
 r_C=[a,CaC],\qquad
 h r_\tau h^{-1},\quad
 r_\tau=[a,\tau a\tau],\quad
 h=a\tau a^{-1}\tau a^{-1}.
\end{equation}
The other eighteen orbifold words vanish freely with the cone squares.
The two ordinary sheets of these words give all forty ordinary
node images, including the four nonempty ones.
It follows directly, in both directions from the entire presentation,
that
\begin{equation}\label{eq:B139-two-normal-generators}
 E_a=
 \bigl((\mathbb Z(a)*Q_f)\rtimes S\bigr)/
       \langle\!\langle r_C,r_\tau\rangle\!\rangle .
\end{equation}
The normal closure is taken in the displayed whole group, so it
retains all signed source words and all fiber conjugates.
Removing the conjugator $h$ here uses that normal closure; it is
not a replacement of one based word by an equal word.

The complete finite-input certificate checks all forty ordinary
words and both factors of each, and all twenty whole orbifold
words. It also checks all signed ordinary and orbifold generator
diagrams, the inverse actions and the selected involution.
The orbifold record does not claim a separate factor check.

\emph{B.139.2. A simultaneous free basis and its ordered action.}
The kernel of $\mathbb Z(a)*Q_f\to Q_f$ is free on
\[
 a_g=gag^{-1}\qquad(g\in Q_f).
\]
Indeed the free-product Schreier scan retains the current $Q_f$
prefix and emits one $a_g^{\pm1}$ for each signed $a$-letter.
Expansion is its inverse. These are therefore simultaneous
free generators. Since the operators in \eqref{eq:B139-operators}
preserve $Q_f$, its section is literal and survives the two
node relations, which lie in the free kernel.

For every signed $s\in S$, define the actual cocycle by
\begin{equation}\label{eq:B139-cocycle}
 \alpha_s(a)=h_sah_s^{-1},\qquad
 h_m=1,\quad h_v=f^{-1},\quad
 h_{st}=\alpha_s(h_t)h_s .
\end{equation}
The inverse rule is
$h_{s^{-1}}=\alpha_s^{-1}(h_s^{-1})$.
The exact action of $Q$ on the free basis is
\[
 (gs)a_j(gs)^{-1}=a_{\,g\alpha_s(j)h_s}.
\]
Its stabilizer at $a_1$ is precisely
$H=\{h_s^{-1}s:s\in S\}$.
The ordered identity \eqref{eq:B139-cocycle} makes this a subgroup,
and projection onto $S$ is an isomorphism, with inverse
$m\mapsto m$, $v\mapsto fv$.
In particular $H\cap Q_f=1$.

The two words in \eqref{eq:B139-two-normal-generators} are
the basis commutators at $\{H,CH\}$ and $\{H,\tau H\}$.
Their entire $Q$-orbit is exactly \eqref{eq:B139-graph}.
The normal closure in the free kernel imposes exactly these
commutations. Conversely the graph presentation and its $Q$-action
satisfy both original words and all actual source equations.
The two inverse maps fix every marked generator, proving
\eqref{eq:B139-split}. No simultaneous independence of arbitrary
peripheral relators is assumed.

\emph{B.139.3. The complete selected edge classification.}
Every element of $S$ fixes $\tau$. Reduced free-product words in
$Q_f$ give $C_{Q_f}(\tau)=\langle\tau\rangle$, and its two cone
factors are not conjugate. Hence
\[
 C_Q(\tau)=\langle\tau\rangle\times S,\qquad
 C_Q(\tau)\cap Q^0=S,
\]
and $C,\tau$ are not conjugate in $Q$.
There is no $\tau$-fixed vertex in $Q/H$: a putative
$q^{-1}\tau q\in H$ projects trivially to $S$, contradicting
injectivity of $H\to S$.

Suppose a $\tau$-exchanged pair is an edge. Write its endpoints
as $rH,rtH$, with $t=C$ or $\tau$.
Then $r^{-1}\tau r=th$ for some $h\in H$.
Projection onto $S$ gives $h=1$.
The $C$ case is impossible by nonconjugacy; the other gives
$r\in C_Q(\tau)$. Thus all such pairs, without an orbit cutoff, are
\[
 \{sH,\tau sH\}\qquad(s\in S).
 \]
Their commuting products are exactly
$\alpha_s(a\delta(a))=\alpha_s(A_s)$, where
$A_s=s_1s_4^{-1}s_1^{-1}$ is the corrected actual contact of B.130.

\emph{B.139.4. Discrete graph fixed elements and the raw source.}
We use the finite-support piling normal form already proved in B.112.
For a graph involution without fixed vertices, a nontrivial fixed
reduced word has an initial signed letter $x^\epsilon$ and also
admits a reduced representative beginning
$\delta(x)^\epsilon$.
The reduced-word interchange property forces these distinct letters
to commute. Both may be moved successively to the front.
Their product is fixed, and removing it leaves a fixed element of
length two less. Induction shows that the fixed subgroup is
generated by the commuting exchanged-pair products.
Every word has finite support, so this argument applies to the
infinite graph; it is not an argument about profinite words.

The six original raw source words project, by free reduction alone, to
\begin{equation}\label{eq:B139-raw-source}
 (s_1,s_2,s_3,s_4,v,m)\longmapsto
       (1,1,1,(aa^*)^{-1},v,m).
\end{equation}
This is checked from the raw B.119/B.123 source marking, not by
substituting node-reduced representatives before projection.
Thus the actual source image is
$I_a\rtimes S$, where
$I_a=\langle\alpha_s(aa^*):s\in S\rangle$.
The preceding classification and induction give
$A(\mathcal G)^\delta=I_a$.
Unique semidirect decomposition and $(Q^0)^\delta=S$
prove the first equality in \eqref{eq:B139-main}.

\emph{B.139.5. Finite graph models refining every actual finite target.}
Put $b_s=h_s^{-1}$; then
$b_{st}=b_s\alpha_s(b_t)$.
Let
\[
 \chi:Q_f\to C_2^2,\qquad
 \chi(C)=(1,0),\quad\chi(\tau)=(0,1),\quad\chi(f)=0.
\]
Both signed source actions preserve $\chi$, and $\chi(h_s)=0$.
Its sum is the original orientation.

Fix any actual finite target $\rho:E_a\to T_0$, augmented by
orientation. Choose a characteristic finite-index normal subgroup
$N_0\leq Q_f$ contained in $\ker(\rho|_{Q_f})\cap\ker\chi$.
Such subgroups are cofinal because $Q_f$ is finitely generated.
Set $F_0=Q_f/N_0$. The actual homomorphism
\[
 S\longrightarrow F_0\rtimes\operatorname{Aut}(F_0),
       \qquad s\longmapsto(b_s,\alpha_s)
\]
records both action and cocycle. Intersect its kernel with
$\ker(\rho|_S)$, obtaining $L_0$, and put
$S_0=S/L_0$, $D_0=F_0\rtimes S_0$.
The image $H_0$ of $H$ is the graph of the descended cocycle,
so $H_0\to S_0$ is injective and $H_0\cap F_0=1$.
Both cone characters descend.

Let $\mathcal G_0$ have vertices $D_0/H_0$ and all $D_0$-translates
of the same two seed edges. They are not loops, by the cone characters
and $H_0\cap F_0=1$. Put
$G_0=A(\mathcal G_0)\rtimes D_0$.
There are actual surjections $E_a\to G_0\to\rho(E_a)$.
For the only well-definedness issue, if $q_1,q_2$ have the same
$D_0/H_0$ label, lift the intervening $H_0$ element to $h\in H$.
Then $q_2^{-1}q_1h^{-1}$ is killed by $\rho|_Q$.
Since $H$ centralizes $a$, the two proposed graph-generator
images agree. Every finite edge lifts to an actual edge, and the
$D_0$ action is respected. This proves factorization of the whole
target, rather than of a selected list of elements.

If $N_1\leq N_0$ is another characteristic finite kernel, choose
$L_1\leq L_0$ also inside its finite action/cocycle kernel.
This gives compatible actual maps $D_1\to D_0$ and
$G_1\to G_0$. No inverse-limit identification of the completion
of the infinite graph group has been made.

\emph{B.139.6. Refining the coarse centralizers.}
The full free-profinite-product factor centralizer theorem gives
$C_{\widehat Q_f}(\tau)=\langle\tau\rangle$
\citep{A_HR1985}.
For the prescribed $F_0$, choose a characteristic refinement $N_1$
such that
\begin{equation}\label{eq:B139-centralizer-refinement}
 \operatorname{im}\bigl(C_{Q_f/N_1}(\tau)\to F_0\bigr)
       \subseteq\langle\tau\rangle.
\end{equation}
Here is the compactness justification. The images of finer finite
centralizers in $F_0$ form a decreasing directed family of finite
sets. Their intersection is the image of
$C_{\widehat Q_f}(\tau)$, by the finite-intersection property for
compatible lifts. This does not require surjective centralizer
transition maps. A decreasing directed family of finite sets
attains its intersection after a finite refinement.

Construct $G_1\to G_0$ using this refinement.
Since all sources fix $\tau$,
$C_{D_i}(\tau)=C_{F_i}(\tau)\rtimes S_i$.
Every element of $C_{D_1}(\tau)$ maps to $\tau^\epsilon s_0$
in $D_0$; if its orientation is even, $\epsilon=0$.
The finite graph edge classification is exactly the argument of
B.139.3: $H_i\to S_i$ excludes fixed vertices, and projection
followed by the separate cone characters excludes exchanged
$C$-edges. Thus $A(\mathcal G_i)^\delta$ is generated by the
norms indexed by $C_{D_i}(\tau)$.
Under $G_1\to G_0$, each such norm maps to an actual source norm
$\alpha_s(A_s)$; the factor $\tau^\epsilon$ only reverses the
order of its commuting pair. The even coarse centralizer maps
to $S_0$. Consequently
\begin{equation}\label{eq:B139-discrete-head-control}
 \operatorname{im}\bigl(C_{G_1^0}(\tau)\to G_0\bigr)
       \subseteq\operatorname{im}(J_a\to G_0).
\end{equation}

\emph{B.139.7. The precise finite-head completion theorem.}
For each finite graph $\mathcal G_i$, the graph group has its
finite Salvetti $K(\pi,1)$, of dimension its maximum clique size.
The full-link special-HNN induction, with the coordinate
retractions onto full subgraphs, proves residual finiteness
and the required full subgroup topologies as in B.112.
A finite extension by $D_i$ remains residually finite:
take the finite intersection of conjugates of a separating
finite-index normal subgroup in the normal graph group.

By \citet[Theorem~3.15 and its complete local dependencies
3.10--3.14]{B_Lorensen2008}, every finite graph group is good
for finite modules. For an arbitrary discrete continuous
$\mathbb F_2[[\widehat G_i]]$-module, the span of finitely many
compact-group orbits is finite. The module is therefore a
filtered union of finite invariant submodules.
The finite cellular resolution makes discrete graph-group
cohomology commute with this union; continuous cochains on
compact domains have finite image, giving the continuous
counterpart. The comparison consequently holds for all these
modules. The finite-extension Lyndon--Hochschild--Serre
comparison transfers it to $G_i$.
The same finite resolution and finite-extension spectral
sequence give finite cohomology in each degree with finite
coefficients, as well as finite virtual $2$-cohomological
dimension. Thus the finite virtual cohomological \emph{type}
and all-module hypotheses are both checked.
The graphs need not be bipartite, and no uniform bound on their
clique dimensions is required.

Apply only the prime-order case of
\citet[Theorem~4.3, Lemma~4.4, Section~4.2 and Lemma~2.2]
{BoggiZalesskii2025} to the actual dense injection
$G_i\hookrightarrow\widehat G_i$ and $\tau$.
It gives
\[
 C_{\widehat G_i}(\tau)=\overline{C_{G_i}(\tau)},\qquad
 C_{\widehat{G_i^0}}(\tau)=\overline{C_{G_i^0}(\tau)}.
\]
The second equality is intersection with the open ordinary
subgroup; density restricts to that open part.
This is not an inference from goodness alone or from a
finite-support induction applied to a completed element.

For $z\in(\widehat{\Gamma_a})^\delta$, map to
$\widehat{G_1^0}$. The displayed centralizer equality and
\eqref{eq:B139-discrete-head-control} show that its image
in the arbitrary finite target $\rho(E_a)$ belongs to
$\rho(J_a)$. This holds for every such target.
Finite-index core refinement in $E_a$ retains all finite
targets of $\Gamma_a$. Closed subgroups of a profinite group
are determined by these inverse finite images, so
$z\in\overline{J_a}$. The reverse inclusion follows because
every actual source generator is fixed.
This proves the completed equality in \eqref{eq:B139-main}.
No residual-finiteness assumption on the original $E_a$
or $\Gamma_0$ was used.

\emph{B.139.8. The exact original remaining kernel.}
Every discrete fixed element in $\Gamma_a$ has an actual
source representative, which is already fixed in $\Gamma_0$.
The fixed projection is therefore onto.
For $z\in\Gamma_0^\delta$, choose such a source lift $h$ of
$\pi(z)$; then $h^{-1}z\in K_{\rm rem}^\delta$.
This gives the first two assertions of
\eqref{eq:B139-kernel-reductions} and their criterion.

Finite-quotient right exactness gives
$\widehat{\Gamma_a}
=\widehat{\Gamma_0}/\overline{K_{\rm rem}}$.
The compact image of $\overline{H_{\rm src}}$ is
$\overline{J_a}$, which is the entire completed fixed subgroup
by \eqref{eq:B139-main}.
The same fixed-lift argument inside the actual completions
gives the last assertion and the closed-kernel criterion.
This particular fixed surjectivity uses the actual source;
it is not a general fixed-lifting assertion for other fillings.
In particular it does not contradict B.135's failure of
fixed lifting from the different inner-star map.
\end{proof}

The full finite word certificates, raw-source verification and
all-target proof records are retained at the following paths.
They check the whole defining input, not a truncated source orbit.
\begin{quote}\footnotesize\raggedright
\path{work_CE/code/compile_one_family_face28.py},
\path{work_CE/code/check_one_family_sources28.py},
\path{work_CE/next_symmetric_nodes/one_family_face28.json},
\path{work_CE/next_symmetric_nodes/one_family_sources28.json},
\path{work_CE/one_family_graph_fixed_source28.txt},
\path{work_CE/one_family_kernel_reduction28.txt},
\path{work_CE/compact_one_family_fixed_source28.txt}.
\end{quote}
